\subsection{Abelian wild inertia and metabelian quotients}

\begin{proposition}[F10: the abelian inertia kernel is finite; Proved]
Let $A,U,X,D=\bigcup_iC_i$ be as in F8. The kernel of
\[
 \pi_1(U)^{\mathrm{ab}}\longrightarrow\pi_1(X)^{\mathrm{ab}}
 \tag{F10.1}
\]
is finite. More explicitly, write $B=(C_i\cdot C_j)$ and set
\[
 k_i=\max\{0,2g_i-2+\nu_i\},\qquad
 t=(-B)^{-1}k.
 \tag{F10.2}
\]
Every $p$-primary rank-one character $\chi$ of $\pi_1(U)$ has
Swan conductor $s_i(\chi)\leq t_i$ at $C_i$.
Choose $e\geq0$ with $p^e>\max_i t_i$. Then the kernel in
(F10.1) is annihilated by $p^e|\det B|$.
\end{proposition}

\paragraph{Status and source.}
The complete argument passed an independent root proof and primary-source
audit; BCE separately passed F10 and F10b. The input used here is the
global refined Swan conductor on regular $F$-finite
schemes in Krishna--Majumder, arXiv:2601.12177v2, 8 November 2026,
Corollary 9.7 p.~56 (Theorem 9.4 pp.~55--56), with the hypotheses
on p.~2. The source is a preprint. Its formulation includes
complete-local resolutions and does not require algebraization
or finite type over a field. No priority assertion is made.

\begin{proof}
A coefficient field identifies $A$ as a quotient of a power-series
ring over perfect $k$, so $A$ and the resolution charts are
$F$-finite. The resolution is an equidimensional noetherian regular
$\mathbf F_p$-scheme with SNC boundary. These are precisely the
hypotheses of the cited global refined-conductor theorem.
All differential forms over $\mathbf F_p$ below may also be viewed
over $k$, since $k$ is perfect and $\Omega^1_{k/\mathbf F_p}=0$.

Fix a character $\chi\in H^1(U,\mathbf Q_p/\mathbf Z_p)$.
It has finite order, say dividing $p^m$. Let
\[
                         S=\sum_i s_i C_i
\]
be its logarithmic Swan-conductor divisor. Each $s_i$ is a
nonnegative integer, defined at the henselian discrete valuation
field at the generic point of $C_i$. If $s_i>0$, the character is
in $\operatorname{fil}^{\log}_S H^1(U,\mathbf Z/p^m)$ and is not
in $\operatorname{fil}^{\log}_{S-C_i}H^1(U,\mathbf Z/p^m)$.
Since $S-C_i\geq\lfloor S/p\rfloor$, Corollary 9.7 of
\cite{KrishnaMajumder2026} gives a nonzero section
\[
 0\ne\operatorname{Rsw}_i(\chi)\in
 H^0\!\left(C_i,
   \Omega_X^1(\log D)|_{C_i}\otimes\mathcal O_X(S)|_{C_i}\right).
 \tag{F10.3}
\]
The map is a levelwise monomorphism, not only an injection after
passing to an unspecified colimit; see Theorem 9.4(1).
Its generic stalk is the local refined Swan conductor by
Corollary 9.6. Thus the nonvanishing in (F10.3) has exactly the
required meaning.

The residue sequence on $C_i$ is
\[
 0\longrightarrow\Omega^1_{C_i/k}(\log\{\text{crossings}\})
 \longrightarrow\Omega_X^1(\log D)|_{C_i}
 \longrightarrow\mathcal O_{C_i}\longrightarrow0.
 \tag{F10.4}
\]
It follows either from the local coordinates $d\log x,dy$ and
$d\log x,d\log y$, or from Lemma 9.8 of the same source.
The two line bundles have degrees $2g_i-2+\nu_i$ and $0$.
Consequently every line subbundle of the middle term has degree
at most $k_i$. A nonzero section in (F10.3), after saturation,
therefore implies
\[
                              S\cdot C_i\geq-k_i.
 \tag{F10.5}
\]
When $s_i=0$, this inequality holds as well, because
$S\cdot C_i=\sum_{j\ne i}s_j(C_j\cdot C_i)\geq0$.
Thus, as a vector inequality, $Bs\geq-k$.

We include the elementary matrix step. Put $C=-B$. This matrix
is symmetric positive definite and has nonpositive off-diagonal
entries. Its inverse is entrywise nonnegative. Indeed, if
$Cx=v\geq0$, write $x=x_+-x_-$ with $x_+,x_-$ nonnegative and
of disjoint support. If $x_-\ne0$, then
\[
 0\leq x_-^Tv=x_-^TCx_+-x_-^TCx_-<0,
\]
because the mixed term is nonpositive and the quadratic term is
positive. This contradiction proves $x\geq0$.
Applying $C^{-1}$ to $Cs\leq k$ now gives $s\leq C^{-1}k=t$.
This proves the uniform conductor bound.

\emph{A fixed $p$-power kills all abelian wild inertia.}
We recall explicitly why multiplication by $p$ reduces the bound
by a factor $p$. For a henselian characteristic-$p$ DVR $R$ with
fraction field $L$, the Brylinski filtration on Witt vectors
$a=(a_{m-1},\ldots,a_0)\in W_m(L)$ is
\[
 \operatorname{fil}_nW_m(L)
   =\{a:p^jv(a_j)\geq-n\text{ for all }j\}.
\]
Its image under Artin--Schreier--Witt is the logarithmic conductor
filtration. Witt restriction drops the last displayed coordinate
and gives
\[
 R:\operatorname{fil}_nW_m(L)
       \longrightarrow\operatorname{fil}_{\lfloor n/p\rfloor}
                         W_{m-1}(L).
 \tag{F10.6}
\]
This is Lemma 3.2, p.~12, of \cite{KrishnaMajumder2026}; the
classical characteristic-$p$ description is also given in
\cite[\S3.2.2--3.2.4, pp.~14--15]{KatoLealSaito2019}.
Under the standard embeddings $\mathbf Z/p^m\hookrightarrow
\mathbf Q_p/\mathbf Z_p$, restriction represents multiplication
by $p$. Hence
\[
 s_i(p^e\chi)\leq\lfloor s_i(\chi)/p^e\rfloor=0.
\]
Conductor zero here means unramified: its Witt representative is
integral, and the Artin--Schreier--Witt equation then defines an
\'etale cover of the DVR. Thus $p^e\chi$ is unramified along every
$C_i$. Purity on the regular $X$ extends it to $X$.

\emph{The prime-to-$p$ part.}
The Kummer/intersection calculation in F5 shows that multiplying
any prime-to-$p$ character by $b=|\det B|$ kills all its inertia
along the $C_i$. Purity again extends the resulting character
to $X$. One can see the calculation directly: extend its torsion
line bundle to $X$, write the divisor of the trivialization of
its $n$th power as $\sum a_i C_i$, and intersect with $C_j$.
Then $Ba\equiv0\pmod n$, so $ba_i\equiv0\pmod n$ by the
adjugate matrix. These residues describe the tame inertia.

Set $Q=\pi_1(U)^{\mathrm{ab}}$ and let $I$ be the kernel in
(F10.1). The map is surjective, since a connected finite \'etale
cover of normal $X$ stays connected on its dense open subset $U$.
The previous two paragraphs show that all characters of $I$
obtained by restriction from $Q$ vanish on $bp^eI$. Finite
characters of the profinite abelian group $Q$ separate its
points, so $bp^eI=0$. By Grothendieck--Murre and Suzuki, $Q$ is
topologically finitely generated. Its subgroup of elements
annihilated by the fixed integer $bp^e$ is finite. Therefore $I$
is finite, as asserted.
\end{proof}

\begin{corollary}[F10b: one fixed cover removes every abelian inertia action; Proved]
There is a finite connected abelian cover $U'\to U$ such that,
for every finite locally constant sheaf on $U$ with abelian
monodromy, its pullback to $U'$ extends to a finite locally constant
sheaf on a fixed SNC resolution $X'$ of the normalization of $A$
in $U'$. The same cover and resolution work for all coefficient
primes and all ranks.
\end{corollary}
\begin{proof}
Choose an open subgroup $Q'\subset Q$ with $Q'\cap I=0$; this
is possible because $I$ is finite and $Q$ is profinite. Let $U'$
be the cover corresponding to the inverse image of $Q'$.
Its normalization $A'$ is finite over $A$, excellent, normal,
complete and local with residue field $k$. Fix an SNC resolution
$X'$ of $A'$. The pulled-back abelian monodromy factors through
$Q'$.

At a height-one exceptional point of $X'$, let $R$ be its DVR.
The morphism $\Spec\operatorname{Frac}(R)\to X$ extends to
$\Spec R\to X$ by properness of $X$ over $A$. Every finite
character of $Q/I$, being a cover of $X$, is unramified after this
pullback. Hence the image in $Q$ of the inertia at $R$ lies in
$I$. It also lies in $Q'$, so it is trivial. Thus the given pulled-back
monodromy is unramified at every such DVR. It is already \'etale
elsewhere. Normalization and regular-scheme purity now extend the
corresponding finite cover to $X'$. This proof uses one fixed
$U'$ and $X'$, independently of the sheaf.
\end{proof}

\begin{proposition}[F10c: all-prime bounds for abelian coefficient monodromy; Proved]
There is a constant $C_A$ such that for every prime $\ell$ and
every finite locally constant $\mathbf F_\ell$-sheaf $M$ on $U$
with abelian monodromy,
\[
                       \dim_{\mathbf F_\ell}H^1(U,M)
                       \leq C_A\operatorname{rank}M.
 \tag{F10.7}
\]
\end{proposition}
\begin{proof}
Use the fixed cover and resolution of F10b. If $\ell=p$, F7
bounds the cohomology after pullback by $C_{X'}\operatorname{rank}M$.
If $\ell\ne p$, every degree-one cocycle after pullback already
factors through the tame quotient. Indeed, it gives an affine
monodromy representation into $M\rtimes\operatorname{im}\rho$.
The original representation $\rho$ is unramified on $X'$, so its
restriction to boundary inertia is trivial. On inertia the cocycle
is a homomorphism to the additive $\ell$-group $M$. It kills wild
pro-$p$ inertia, and the affine cover is therefore tame.
Consequently
\[
 H^1(U',M)=H^1(\pi_1^t(U';X',D'),M),
\]
whose dimension is at most $d_{X'}\operatorname{rank}M$ by F9
and the values of a cocycle on generators.

Let $F=\operatorname{Gal}(U'/U)$. Inflation--restriction gives
\[
 0\to H^1(F,M^{\pi_1(U')})\to H^1(U,M)
          \to H^1(U',M)^F.
\]
The first term has dimension at most $d(F)\operatorname{rank}M$,
with no assumption that $|F|$ is prime to $\ell$.
Thus one can take
$C_A=d(F)+\max\{C_{X'},d_{X'}\}$.
\end{proof}

\begin{theorem}[F11: all metabelian local quotients are finitely generated; Proved]
Every metabelian quotient $G$ of $\pi_1(U)$ is topologically
finitely generated, with a bound depending only on $A$.
Both tame and wild normal factors are allowed.
\end{theorem}
\begin{proof}
Let $d_0=d(\pi_1(U)^{\mathrm{ab}})$ and let $C_A$ be as in
F10c. Put $D=\max\{d_0,\lceil C_A\rceil+1\}$.
We show that every finite metabelian quotient $F$ is generated
by $D$ elements. If $F'=1$, the bound $d_0$ applies. Otherwise
choose a minimal nonzero $F$-normal subgroup $M\subset F'$.
It is an elementary abelian $\ell$-group for some prime $\ell$,
and is simple as an $\mathbf F_\ell[F]$-module. Since $F'$ is
abelian, conjugation on $M$ factors through the abelian group
$F/F'$. Inflation from $F$ to $\pi_1(U)$ and F10c give
\[
                  \dim H^1(F,M)\leq C_A\dim M.
 \tag{F11.1}
\]
Inductively choose $D$ generators of $F/M$ and consider their
$|M|^D$ possible lifts. A nongenerating lift spans a proper
subgroup $J$ with $JM=F$. The subgroup $J\cap M$ is $F$-stable:
$J$ preserves it, and $M$ centralizes it. Minimality therefore
forces $J\cap M=0$, so $J$ is a complement. If there is no
complement, all lifts generate.

In the split case write $H=F/M$. Complements are indexed by
$Z^1(H,M)$. Since $M$ acts trivially on itself, splitting gives
\[
 H^1(F,M)\simeq H^1(H,M)\oplus\End_H(M).
\]
The endomorphism ring contains the scalar field $\mathbf F_\ell$.
Moreover $|B^1(H,M)|\leq|M|$. Hence (F11.1) implies
\[
 |Z^1(H,M)|\leq |M|^{C_A+1}/\ell<|M|^D.
\]
Each complement contains at most one lift of the chosen generating
tuple of $H$, so some lift generates $F$. This proves the bound
by induction. Finally the compactness of generating-tuple conditions
on $G^D$ gives $D$ topological generators of $G$.
\end{proof}

\paragraph{Scope remaining after these results.}
F11 closes derived length two, including the metabelian
abstract obstructions of F3 in every residue characteristic.
It does not treat arbitrary solvable quotients of larger derived
length: their minimal normal modules can have nonabelian
coefficient monodromy. Nor does it settle arbitrary nonsolvable
quotients. These are genuine remaining gates.
